Considera el sistema
\[
\left.\begin{aligned}
x+2y-z&=3\\
2x+y+z&=3\\
x-y+2z&=0
\end{aligned}\right\}
\]

\begin{apartats}

\apartat[1,25]{0,75}
Escriu-lo en forma matricial i classifica'l amb el teorema de Rouché-Fröbenius.

\begin{solucio}
$\begin{pmatrix}1&2&-1\\2&1&1\\1&-1&2\end{pmatrix}\begin{pmatrix}x\\y\\z\end{pmatrix}=\begin{pmatrix}3\\3\\0\end{pmatrix}$.
Sigui $A$ la matriu de coeficients i $A^*$ l'ampliada.\\
$|A|=2+2+2+1-8+1=0$ i $\begin{vmatrix}1&2\\2&1\end{vmatrix}=-3\ne0$: $\operatorname{rang}A=2$.\\
A $A^*$, la tercera fila és la segona menys la primera, també al terme independent ($0=3-3$): tots els
menors d'ordre 3 són nuls i $\operatorname{rang}A^*=2$.\\
$\operatorname{rang}A=\operatorname{rang}A^*=2<3$: és compatible indeterminat, amb un paràmetre.
\end{solucio}

\begin{tria}{rouche-tasca}
\itemtria{resol-indeterminat}{1}{1,25}
Resol el sistema i dona'n dues solucions particulars.

\begin{solucio}
N'hi ha prou amb les dues primeres equacions. Amb $z=\lambda$:
$\left.\begin{aligned}x+2y&=3+\lambda\\2x+y&=3-\lambda\end{aligned}\right\}$. Restant a la segona el doble
de la primera, $-3y=-3-3\lambda$, $y=1+\lambda$, i $x=3+\lambda-2y=1-\lambda$.\\
Les solucions són $(x,y,z)=(1-\lambda,\ 1+\lambda,\ \lambda)$. Per exemple, amb $\lambda=0$, $(1,1,0)$, i
amb $\lambda=1$, $(0,2,1)$.
\end{solucio}

\itemtria{mes-equacions}{1}{1,25}
Estudia amb el teorema de Rouché-Fröbenius aquest altre sistema, de tres equacions i dues incògnites, i
resol-lo si és compatible:
$\left.\begin{aligned}x+2y&=5\\3x-y&=1\\2x+y&=4\end{aligned}\right\}$

\begin{solucio}
$\begin{vmatrix}1&2\\3&-1\end{vmatrix}=-7\ne0$: el rang de la matriu de coeficients és 2, el màxim.
L'ampliada és quadrada: $\begin{vmatrix}1&2&5\\3&-1&1\\2&1&4\end{vmatrix}=-4+4+15+10-1-24=0$, i el seu
rang també és 2.\\
$\operatorname{rang}A=\operatorname{rang}A^*=2=$ nombre d'incògnites: és compatible determinat. De les
dues primeres, $x=1$ i $y=2$, que també compleixen la tercera: $2+2=4$.
\end{solucio}
\end{tria}

\begin{nomesllarg}
\apartat{0,75}
Classifica el sistema $\left.\begin{aligned}x-y+z&=1\\-2x+2y-2z&=3\end{aligned}\right\}$. Canvia'n el
terme independent de la segona equació perquè sigui compatible. De quants paràmetres dependran llavors
les solucions?

\begin{solucio}
Les dues files de coeficients són proporcionals: $\operatorname{rang}A=1$. A l'ampliada,
$\begin{vmatrix}1&1\\-2&3\end{vmatrix}=5\ne0$: $\operatorname{rang}A^*=2$. És incompatible.\\
Amb terme independent $-2$, la segona equació és la primera multiplicada per $-2$:
$\operatorname{rang}A=\operatorname{rang}A^*=1<3$. És compatible indeterminat, i les solucions depenen
de $3-1=2$ paràmetres.
\end{solucio}
\end{nomesllarg}

\end{apartats}
